% Challenge dossier: VI.05.C1
\subsection*{Challenge VI.05.C1 --- The component-intersection graph}
\label{chal:vi-05-c1}

\paragraph{Problem.}
\label{prob:vi05-c1}
Let
\[
X=X_1\cup\cdots\cup X_r
\]
be the decomposition of an affine algebraic set into irreducible components.  Form the graph
whose vertices are the \(X_i\), joining \(X_i\) and \(X_j\) when
\[
X_i\cap X_j\neq\varnothing.
\]
Prove that the connected components of this graph correspond exactly to the connected
components of \(X\).

\ifdefined\FullProblemDossiers

\paragraph{Why this problem matters.}
This challenge turns the component-intersection picture into a precise theorem:
topological connectedness of the whole algebraic set is exactly graph connectedness of its
irreducible components.

\paragraph{Useful ideas before starting.}
\begin{itemize}
\item irreducible components are connected;
\item a union of connected subsets linked by nonempty intersections is connected;
\item finite unions of closed subsets are closed.
\end{itemize}

\paragraph{Guided hints.}
\textbf{Hint 1.} For one connected component of the graph, take the union of the
corresponding irreducible components.

\textbf{Hint 2.} Use paths in the graph to prove that this union is connected.

\textbf{Hint 3.} Distinct graph-components have no intersections between their corresponding
unions, so these unions form a separation of \(X\).

\paragraph{Solution.}
Let the connected components of the graph be
\[
\Gamma_1,\dots,\Gamma_m.
\]
For each graph-component define
\[
Y_\alpha
=
\bigcup_{X_i\in\Gamma_\alpha}X_i.
\]
This is a finite union of closed subsets, hence closed.

Within a fixed \(\Gamma_\alpha\), any two vertices are joined by a chain
\[
X_{i_0},X_{i_1},\dots,X_{i_r}
\]
with consecutive intersections nonempty.  Since every \(X_i\) is irreducible and therefore
connected, repeated use of the fact that a union of intersecting connected subsets is
connected shows that \(Y_\alpha\) is connected.

If \(\alpha\neq\beta\), no component in \(\Gamma_\alpha\) meets a component in
\(\Gamma_\beta\), so
\[
Y_\alpha\cap Y_\beta=\varnothing.
\]
Also,
\[
X=Y_1\sqcup\cdots\sqcup Y_m.
\]
Because each \(Y_\alpha\) is closed and its complement is the finite union of the other
\(Y_\beta\), each \(Y_\alpha\) is also open.  Therefore the \(Y_\alpha\) are exactly the
connected components of \(X\).

\paragraph{What happened mathematically.}
The theorem compresses a topological question into finite combinatorics once the
irreducible components are known.  In practical examples this is often the fastest way to
decide connectedness after decomposition.

\paragraph{Extensions.}
\begin{enumerate}
\item Apply the criterion to a chain, a cycle, and a star of affine lines.
\item Construct an algebraic set whose component graph is any prescribed finite graph.
\item Relate the graph components to primitive idempotent decompositions of the coordinate
ring when the base field is algebraically closed.
\end{enumerate}

\paragraph{Provenance.}
New canonical challenge synthesized from the chapter's component-intersection criterion.

\fi

